\section{Line integrals}

To start inspecting new types of integrals, we have to get into differential forms. We start with the simplest type. In what comes, we identify the covector space $\mathbb{R}^{n*}$ with the set of row vectors (they're canonically isomorphic anyway so we might as well). And I for the first time in history start paying attention to the difference between row and column vectors.

\begin{definition}[$1$-forms]\label{def:1-forms}
	Let $U \subset \mathbb{R}^{n}$ be open. A map assigning to each point $x \in U$ a covector $x \mapsto (\alpha_{1}(x), \dots, \alpha_{n}(x)) = \alpha_{x} \in \mathbb{R}^{n*}$ is called a $1$-form. We say that a $1$-form is $C^{k}$ if every $\alpha_{i}: U \to \mathbb{R}$ is $C^{k}.$
\end{definition}
\newpage
\begin{eg}[Exterior derivative]
	For $f \in C^{k}(U, \mathbb{R})$, notice that $x \mapsto Df(x)$ is a $1$-form. We will talk about this in more detail maybe later, but it makes sense to call $f$ a $0$-form. Then the map $df: U \to \mathbb{R}^{n*}: x \mapsto Df(x)$ turns the $0$-form $f$ into a $1$-form $df$.
\end{eg}

\begin{definition}[Line integral]\label{def:line-int}
	If $\alpha: U \to \mathbb{R}^{n*}$ is a class $C^{1}$ $1$-form in $U \subset \mathbb{R}^{n}$ open and $\gamma : [a,b] \to U$ a $C^{1}$ curve, then we define the line integral as
	\begin{align*}
		\int_{\gamma } \alpha := \int_{a}^{b} \alpha_{\gamma (t)}(\gamma '(t)) \, dt = \int_{a}^{b} \alpha_{\gamma (t)} \cdot \gamma '(t) \, dt.
	\end{align*}
\end{definition}

\begin{lemma}[Invariance under positive reparametrisation]\label{lem:inv-pos-rep}
	Let $U \subset \mathbb{R}^{n}$ open, $\gamma: [a,b] \to U \subset \mathbb{R}^{n}$ a curve and $\psi: [c,d] \to [a,b]$ a $C^{1}$ diffeomorphism such that $\psi' > 0$. Then
	\begin{align*}
		\int_{\gamma } \alpha = \int_{\gamma \circ \psi} \alpha .
	\end{align*}
\end{lemma}
\begin{proof}
	We can substitute $t = \psi(s)$ and we get
	\begin{align*}
		\int_{\gamma \circ \psi} \alpha = \int_{c}^{d} \alpha_{\gamma \circ \psi(t)} ((\gamma \circ \psi)'(s)) \, ds = \int_{c}^{d} \alpha_{\gamma \circ \psi(s)}(\gamma '(\psi(s))) \psi'(s) \, ds = \int_{a}^{b} \alpha_{\gamma (t)} \gamma '(t) \, dt = \int_{\gamma } \alpha.
	\end{align*}
\end{proof}

\begin{definition}[Conservative vector field / Exact $1$-form]\label{def:con-vec-field}
	Let $U \subset \mathbb{R}^{n}$ be connected and open. A $1$-form $\alpha : U \to \mathbb{R}^{n*}$ class $C^{1}$ is \textit{exact} if $\exists f \in C^{2}(U)$ such that $\alpha = df$.
\end{definition}

\begin{remark}[Parallels to physics and familiar notation]
	Consider a vector field $F$ and the associated $1$-form $\alpha = F^\top$. Then
	\begin{align*}
		\int_{\gamma } \alpha = \int_{\gamma } F^\top = \int_{a}^{b} F^\top (\gamma (t)) \gamma '(t) \, dt = \int_{a}^{b} F(\gamma (t)) \cdot \gamma '(t) \, dt = \int_{\gamma } F(s) \cdot ds
	\end{align*}
	is just the usual "work along a path" from physics.

	We say a vector field $F \in C^{1}(U, \mathbb{R}^{n})$ is \textit{conservative} if $\exists f: U \to \mathbb{R}$ $C^{2}$ such that $\nabla f = F$, which is the analogue to "exact".
\end{remark}


\begin{lemma}[Path independence]\label{lem:path-indep}
	Let $\alpha : U \to \mathbb{R}^{n*}$ be an exact $1$-form such that $\alpha = df$ for $f \in C^{2}(U)$. Then for any $C^{1}$ curve $\gamma: [a,b] \to U$,
	\begin{align*}
		\int_{\gamma } \alpha = f(\gamma (b)) - f(\gamma (a)).
	\end{align*}
\end{lemma}
\begin{proof}
	We compute using the chain rule
	\begin{align*}
		\int_{\gamma } \alpha  = \int_{a}^{b} \alpha_{\gamma (t)}(\gamma '(t)) \, dt = \int_{a}^{b} df_{\gamma (t)} (\gamma '(t)) \, dt = \int_{a}^{b} (f \circ \gamma )' \, dt = f(\gamma (b)) - f(\gamma (a)).
	\end{align*}
\end{proof}

\begin{lemma}[Integrability conditions]\label{lem:int-cond}
	Let $U \subset \mathbb{R}^{n}$ open, $\alpha : U \to \mathbb{R}^{n*}$ be exact ($\alpha = df$) and $C^{1}$. Then $\partial_{j} \alpha_{i} = \partial_{i} \alpha_{j}$.
\end{lemma}
\begin{proof}
	Schwarz Theorem \ref{thm:schwarz} applies: $\partial_{j} \alpha_{i} = \partial_{j} \partial_{i} f = \partial_{i} \partial_{j} f = \partial_{i} \alpha_{j}$.
\end{proof}

\begin{lemma}[Converse on convex domains]\label{lem:conv-conv-dom}
	Assume $U \subset \mathbb{R}^{n}$ is open convex and let $\alpha: U \to \mathbb{R}^{n*}$ be a $1$-form that satisfies the integrability conditions, i.e. $\partial_{i} \alpha_{j} = \partial_{j} \alpha_{i}$. Then $\exists f: U \to \mathbb{R}$ class $C^{2}$ such that $df = \alpha$.
\end{lemma}
\begin{proof}
	% Assume first the domain is a box in $\mathbb{R}^{2}$, so $U = (a,b) \times (c,d)$.
	% Fix $(x_{0}, y_{0}) \in U$ and define
	% \begin{align*}
	%     f(x,y) := \int_{x_{0}}^{x} \alpha_{(t, y_{0})}((1,0)^\top )  \, dt = 
	%     \int_{y_{0}}^{y} \alpha_{(x, s)} ((0,1)^\top ) \, dt.
	% \end{align*}
	% Since $\alpha_{(x,y)} = (\alpha_{1}(x,y), \alpha_{2}(x,y))$, this equals
	% \begin{align*}
	%     f(x,y) = \int_{x_{0}}^{x} \alpha_{1}(t, y_{0})  \, dt + \int_{y_{0}}^{y} \alpha_{2}(x, s)  \, ds.
	% \end{align*}  
	% We claim that $df = \alpha $. By \ref{thm:diff-und-int} and the FTC,
	% \begin{align*}
	%     \partial_{1} f&= \alpha_{1}(x, y) + \int_{y_{0}}^{y} \partial_{1} \alpha_{2}(x,s)  \, ds \\ 
	%     \partial_{2} f&= \int_{x_{0}}^{x} \partial_{2} \alpha_{1} (t, y_{0})  \, dt + \alpha_{2}(x,y).
	% \end{align*}
	% Using the integrability conditions and the FTC,
	% \begin{align*}
	%     \int_{y_{0}}^{y} \partial_{1}\alpha_{2}(x,s)  \, ds &= \int_{y_{0}}^{y} \partial_{2} \alpha_{1}(x, s)  \, ds = \alpha_{1}(x, y) - \alpha_{1}(x, y_{0}) \\ 
	%     \int_{x_{0}}^{x} \partial_{2} \alpha_{1} (t, y_{0})  \, dt = \int_{x_{0}}^{x} \partial_{1} \alpha_{2} (t, y_{0})  \, dt = 0.
	% \end{align*}
	% The rest cancels out, giving the statement. 


	Assume without loss of generality that $0 \in U$ (otherwise translate). We use the path integral of $\alpha$ along the straight line from $0$ to $x \in U$ to define a function $f: U \rightarrow \mathbb{R}$ by
	$$
		f(x):= \int_{\gamma } \alpha = \int_0^1 \alpha(t x) \cdot x \, dt
	$$
	for $x \in U$ and $\gamma(t) = tx$. Let us prove that $\alpha = df$.

	Indeed, fix $j \in\{1, \ldots, n\}$ and consider, as a preparation for the computation of $\partial_j f$,
	\begin{align*}
		\partial_j(\alpha(tx) \cdot x) & =\partial_j\left(\sum_{k=1}^n \alpha_k(t x) x_k\right)
		=\sum_{k=1}^n \partial_j \alpha_k(t x) t x_k+\alpha _j(t x)
		\\ &=  t\sum_{k=1}^n \partial_k \alpha _j(t x) x_k+\alpha _j(t x)
		=\partial_t\left(t \alpha _j(t x)\right)
	\end{align*}
	where we used the integrability conditions.
	
    Hence, using Theorem \ref{thm:diff-und-int}, for $x \in U$ we can compute $\partial_j f(x)$ as
	$$
		\partial_j f(x)=\partial_j \int_0^1 \alpha(t x) \cdot x d t=\int_0^1 \partial_t\left(t \alpha _j(t x)\right) d t=\left[t \alpha _j(t x)\right]_{t=0}^{t=1}=\alpha _j(x) .
	$$
	Thus, $\alpha =d f$, and $f$ is $C^2$.
\end{proof}

